\chapter{Alpha As A Bracket}

\section{The internal formula}

The internal constant \(\alpha_c\) is a computation-owned ratio in the model of one electron reading.  Its consumed
inputs are a target integer \(t>0\), a natural-unit integer \(R>0\), and a positive rational distance
\[
                         d=p/q,\qquad p,q\in\mathbb N_{>0},\quad p>q.
\]
The dependency graph is
\[
(p,q)\longrightarrow\frac{p-q}{p},\qquad
t\longrightarrow\frac{t}{R},\qquad
\left(\frac{p-q}{p},\frac{t}{R},R\right)
\longrightarrow
\alpha_c=\frac{t(p-q)}{pR^2}.
\]
Thus
\[
                         \alpha_c^{-1}=\frac{pR^2}{t(p-q)}.
\]
Every denominator condition is part of the domain.  The formula uses the distance numerator and denominator, target,
and natural-unit scale.  Other report fields---labels, flags, calibration-tape lengths, field-shaped annotations, or
stored coefficients not multiplied into this expression---are unused metadata for this value.

The second-variation ledger supplies the integers \(C=18\) and \(t=5\) for a crossing condition
\[
                         d_\ast^2=C/t=18/5.
\]
This condition determines a model-owned distance question.  Numerical routes approximate its positive root and then
evaluate the exact rational formula above at their rational distance reports.  A stored second-difference value earns
no role unless a dependency edge shows that it supplies \(C\) or \(t\).

For a rational report \(p/q\), the scaled inverse at \(M>0\) is computed by exact floor division:
\[
 I_M=\left\lfloor
 \frac{pR^2M}{t(p-q)}
 \right\rfloor,\qquad
 pR^2M=I_M\,t(p-q)+\rho_M,
\quad0\leq\rho_M<t(p-q).
\]
Retaining numerator, denominator, \(M,I_M,\rho_M\) makes the display loss replayable.  A zero target, \(p=q\), or
zero scale denominator rejects the inverse slice rather than producing a sentinel number.

Two numerical routes are reported separately.  The direct route evaluates the positive crossing in one declared
approximation and then slices the inverse ratio.  The stepped route iterates rational distances
\[
                         d_0\to d_1\to\cdots\to d_k
\]
through declared intermediate updates, retaining every residual, accepted and rejected step, denominator growth,
and stop before applying the same inverse-ratio slice.

At scale \(10^{18}\), the direct and stepped reports are exactly
\[
\begin{aligned}
I_d&=137011290548979455469,\\
I_s&=137011290753751157155,\\
I_d&<I_s,\\
I_s-I_d&=204771701686.
\end{aligned}
\]
These are unequal integers from different processes.  Their difference is read exactly.  No claim of nearness,
equality, containment, or physical agreement follows.

\paragraph{Worked example 1: replaying one rational distance (computation).}
Given \(p/q>1\), compute \(p-q\), form numerator \(t(p-q)\), denominator \(pR^2\), normalize by a positive greatest
common divisor if requested, and perform scaled inverse floor division.  Replay verifies every multiplication,
denominator condition, quotient, and remainder.  The distance \(q\) cancels from the simplified ratio only after the
displayed exact algebra; it remains provenance for the distance report.

\paragraph{Counterexample 2: stored fields are not dependencies (computation).}
A report can contain a Boolean called closure and a coefficient called coupling while the arithmetic consumes neither.
Their presence does not affect \(\alpha_c\).  Altering an unused field while the dependency graph and consumed values
remain fixed leaves this numerical formula unchanged; any semantic claim attached to that field requires a separate
certificate.

\paragraph{Worked example 3: stepped iteration versus final slice (computation).}
The stepped route produces a trace of rational states and resource costs.  Its last state is then projected to
\((I_s,\rho_s)\).  Keeping only \(I_s\) erases the iterative route, rejected steps, exact rational, and floor residue.
The slice can be compared numerically with the direct slice, but it cannot establish process equivalence.

\paragraph{Vacuum draw: Gosset projection exposes what a slice forgets (computation).}
The Gosset experiment sharpens that distinction without supplying another alpha route.  Its finite record domain is
\(H=\mathbb Z^2\), with \(x=(x_u,x_\perp)\), and the named hypothesis direction is \(u=(1,0)\).  The performed
transform is the pair
\[
 T_G(x)=\bigl(\operatorname{proj}_u x,\operatorname{res}_u x\bigr)
       =\bigl((x_u,0),(0,x_\perp)\bigr).
\]
It demonstrates three exact properties over this domain: projection is idempotent, the residual is orthogonal to
\(u\) under the integer dot product, and projection plus residual reconstructs the record coordinatewise.  The
exported experiment then applies the deliberately lossy projection
\(\pi_G(x)=\langle x,u\rangle=x_u\) to a pair of records.  Its canonical inputs \((3,5)\) and \((3,9)\) are distinct
but both return the same shadow \(3\); replacing the second record by \((1,0)\) makes the run return false.  There is
no iterative search: equality of shadows is decided and the run stops, while the failed control exposes the boundary.

The canonical decompositions are \(T_G(3,5)=((3,0),(0,5))\) and
\(T_G(3,9)=((3,0),(0,9))\).  The residual dot product with \(u\) is zero in both cases, but the residuals are unequal.
If \texttt{residual} were replaced by the original record, the check would read
\(\langle(3,5),u\rangle=3\ne0\) and fail.  Thus orthogonality depends on subtracting the projection; it is not a
consequence of calling the second coordinate residual.  Because \(u=(1,0)\) is fixed, changing the axis or dot
product defines another projection.

The commuting statement is exact but narrow:
\[
                         \pi_G(\operatorname{proj}_u x)=\pi_G(x).
\]
The reconstruction statement needs both outputs of \(T_G\); it does not hold from \(\pi_G(x)\) alone.  This is the
same structural warning as the final scaled-integer slice: equality after a projection can coexist with inequality
of sources.  It does not imply that the two alpha routes share an error, distribution, or process.

The source description mentions a normalized statistic, but the executable computes no norm or division.  Its
evidence is the exact projection/residual decomposition and witnessed fiber, not a Student distribution or p-value.

Richardson extrapolation is not licensed merely because direct and stepped routes yield two values.  They are not
declared nested mesh levels of one asymptotic family, and no common error expansion or independently stabilized order
is supplied.  Choosing an order from these two reports to manufacture another value is rejected.

The dependency audit topologically replays raw inputs, crossing data, rational distance construction, exact ratio
products, normalization, scaled-floor quotient and remainder, route iterations, projections, and comparison rows.
It records generated versus consumed fields and rejects orphan, cyclic, or type-incompatible edges.

Budgets distinguish crossing evaluation, direct approximation, stepped iterations, rational normalization,
multiplication and division, denominator and bit growth, scale conversion, storage, comparison, and replay.  Stops
include \textsc{complete}, \textsc{invalid-target}, \textsc{invalid-unit-scale}, \textsc{distance-domain},
\textsc{zero-inverse-denominator}, \textsc{iteration-budget}, \textsc{arithmetic-budget},
\textsc{bit-growth}, \textsc{overflow}, \textsc{nonconvergent}, \textsc{dependency-missing},
\textsc{unused-field-claim}, \textsc{projection-loss-unbounded}, and \textsc{replay-failure}.

A certificate records \(C,t,R,p,q\), crossing and distance domains, route identity, every direct or stepped state,
accepted and rejected updates, consumed inputs, unused stored fields, exact products, normalized ratio, scale,
quotient and remainder, bit lengths, projections and erasures, budgets, stops, checked scope, unresolved frontier,
releases, dependency lock, snapshot hash, configuration, protocol, run identity, platform, deterministic inputs, and
declared nondeterministic inputs.

Comparison reports exact equality, resolution equality under a declared quantizer, and uncertainty-region
compatibility as independent pass/fail/unresolved rows.  Here the exact scaled-integer row fails because \(I_d\ne I_s\).
Any resolution row states only its quantizer result; no uncertainty row exists without a declared uncertainty model.
Same-process additionally requires matching inputs, update law, states, traces, costs, stops, configuration, and
provenance; direct and stepped routes do not satisfy that identity.

Governance fixes the internal formula, domains, consumed dependency graph, crossing data, distance representation,
direct and stepped algorithms, arithmetic and normalization, scale and projection, budgets, stop priority, comparison
relations, and certificate grammar.  A changed field begins a new report.

\(\alpha_c\) is the compiler structure constant and model-owned electron reading.  It is not the laboratory
fine-structure value \(137.036\), and it is not a physical coupling, self-energy, field strength, or apparatus
observation.  The model cannot bridge its name to the world's answer.  A physical comparison would require a
covariant apparatus action, calibrated observation model, uncertainty, interaction trace, back-action, and
provenance.

The action ledger is exact.  Formula, domains, crossing inputs, distance representation, routes, arithmetic, scale,
dependency graph, comparisons, budgets, and certificate grammar are \emph{posited}.  Exact ratio identities,
floor-division identities, dependency reachability, iteration replay, cross-route arithmetic comparisons, and
certificate implications are \emph{demonstrated}.  Inputs, route states, consumed fields, unused fields, products,
scaled integers, difference, remainders, resources, stops, and provenance are \emph{read}.  Equality or nearness of
the unequal reports, process identity, Richardson extrapolation without an asymptotic family, physical meaning,
external alpha, and every bridge from name resemblance remain \emph{reserved}.
\section{The count-to-three bracket}

The crossing is the positive quadratic surd
\[
                         d_\ast=\sqrt{\frac{18}{5}}=\frac{\sqrt{90}}5.
\]
Its simple continued fraction is
\[
                         d_\ast=[1;\overline{1,8,1,2}],
\]
whose convergents begin
\[
                         \frac11,\quad\frac21,\quad\frac{17}{9},
\quad\frac{19}{10},\quad\frac{55}{29},\quad\frac{74}{39},\ldots.
\]
Each partial quotient is obtained by floor division followed by reciprocal of the positive fractional remainder.
The repeating block describes this surd; it is not inferred from decimal formatting.

The recurrence
\[
p_k=a_kp_{k-1}+p_{k-2},\qquad
q_k=a_kq_{k-1}+q_{k-2}
\]
with standard initial pairs generates exact integer convergents.  Replay retains every partial quotient, product,
sum, numerator, denominator, determinant
\[
                         p_kq_{k-1}-p_{k-1}q_k=(-1)^{k-1},
\]
and comparison with the quadratic equation.

The parameter \(K=3\) means three whole partial quotients:
\[
                         [1;1,8]=17/9.
\]
The adjacent previous convergent is \(2/1\).  Exact comparisons give
\[
\left(\frac{17}{9}\right)^2=\frac{289}{81}<\frac{18}{5}<4
=\left(\frac21\right)^2,
\]
since \(1445<1458\) and \(18<20\).  Positivity therefore gives
\[
                         \frac{17}{9}<d_\ast<2.
\]

The internal inverse-ratio map from the preceding section is
\[
                         I(d)=\frac{R^2}{t}\frac{d}{d-1},
\qquad R=18,\quad t=5,\quad d>1.
\]
It is strictly decreasing on this domain.  Evaluating the two rational walls exactly yields
\[
\begin{aligned}
I(2)&=\frac{18^2}{5}\frac{2}{1}=\frac{648}{5}=129.6,\\
I(17/9)&=\frac{18^2}{5}\frac{17}{8}
=\frac{5508}{40}=\frac{1377}{10}=137.7.
\end{aligned}
\]
Thus the \(K=3\) owned jar is the open interval
\[
                         \boxed{(129.6,\ 137.7)}.
\]
Open endpoints mean the exact surd crossing is strictly between the rational distance walls and its mapped internal
value is strictly between their inverse-ratio images.  The endpoints are owned diagnostics computed from adjacent
convergents; neither endpoint is reported as the internal point value.

\paragraph{Worked example 1: exact wall ordering (computation).}
The lower inverse wall is \(648/5\), the upper \(1377/10\), and
\[
                         \frac{648}{5}<\frac{1377}{10}
\quad\Longleftrightarrow\quad6480<6885.
\]
This orders the report without floating arithmetic.  Formatting the rationals as \(129.6\) and \(137.7\) is a
lossless decimal representation here because both denominators divide \(10\).

\paragraph{Worked example 2: one more partial quotient is a new report (computation).}
At \(K=4\), the next convergent is \(19/10\) and the adjacent pair changes.  The resulting wall computation belongs
to a different resolution-depth report.  It cannot be inserted into the \(K=3\) certificate while retaining the
same stop and ownership.

The finite algorithm iterates Euclidean floor and reciprocal steps, slices the convergent stream at \(K\), selects
the adjacent pair, checks straddling against \(5d^2-18\), applies the exact inverse-ratio map to each wall, reverses
the wall order because the map decreases, and emits an open interval.  Every domain, orientation, and ordering check
is recorded.

The three-point Gauss--Legendre rule is a separate marked diagnostic.  It samples an integrand at three nodes with
weights
\[
                         5/9,\quad8/9,\quad5/9.
\]
Its three counts quadrature nodes, not continued-fraction partial quotients.  It neither generates the wall pair nor
supplies the interval width, and the jar is not a quadrature remainder.  Results from that diagnostic require their
own integrand, interval transform, node approximation, arithmetic, and error record.

The older dyadic-grid interval is superseded and is not an alternative owned bracket.  No grid seed, denominator
doubling, or midpoint bisection participates in this continued-fraction certificate.

\paragraph{Counterexample 3: two threes do not identify algorithms (computation).}
A \(K=3\) continued-fraction report and a three-node quadrature report both contain the integer \(3\).  Their inputs,
iterations, data structures, outputs, error models, and stops differ.  Equal counts do not establish same-number or
same-process.

\paragraph{Counterexample 4: a wall is not a selected value (computation).}
Returning \(129.6\), \(137.7\), their midpoint, or a formatted internal diagnostic as the jar's point would add a
selection rule absent from the report.  The owned object is the ordered open interval and its construction trace.

Budgets distinguish partial quotients, Euclidean divisions, numerator and denominator growth, quadratic comparisons,
wall-map arithmetic, scale formatting, storage, and replay.  Stops include \textsc{complete},
\textsc{invalid-radicand}, \textsc{nonpositive-remainder}, \textsc{partial-quotient-budget},
\textsc{insufficient-convergents}, \textsc{bit-growth}, \textsc{overflow}, \textsc{straddling-failure},
\textsc{wall-domain}, \textsc{ordering-failure}, \textsc{mapping-failure}, and \textsc{replay-failure}.

A certificate records the exact radicand, continued-fraction convention, partial quotients, convergent recurrence and
states, determinant checks, \(K\), adjacent-wall selection, exact quadratic comparisons, inverse-ratio map and domain,
wall products and reductions, open-endpoint status, projections and formatting, budgets, stops, checked scope,
unresolved frontier, releases, dependency lock, snapshot hash, configuration, protocol, run identity, platform,
deterministic inputs, and declared nondeterministic inputs.

Comparison reports exact equality, resolution equality under a declared quantizer, and uncertainty compatibility as
separate rows.  The jar itself is an owned interval report, not an uncertainty region unless an uncertainty model
explicitly makes it one.  Same-process requires identical continued-fraction convention, \(K\), recurrence, trace,
wall map, arithmetic, budgets, configuration, and provenance.

Governance fixes the crossing equation, positive-root domain, continued-fraction convention, \(K\), adjacent-wall
policy, straddling test, internal inverse map, arithmetic and formatting, budgets, stop priority, comparison relations,
and certificate grammar.  A changed field begins a new report.

The jar belongs to internal \(\alpha_c\).  The laboratory number \(137.036\) is not inserted, compared by proximity,
or claimed to lie inside it.  No statement that a true coupling lies below a floor is made.  A physical comparison
would require the world-facing apparatus, observation model, calibration, uncertainty, interaction trace, back-action,
and provenance that the computation cannot supply from its own names.

The action ledger is exact.  Crossing, continued-fraction convention, \(K\), wall policy, inverse map, algorithms,
arithmetic, budgets, and certificate grammar are \emph{posited}.  Surd identities, recurrence and determinant
identities, exact straddling, monotonic wall mapping, endpoint derivations, ordering, and certificate implications are
\emph{demonstrated}.  Partial quotients, convergents, wall values, comparisons, resources, stops, and provenance are
\emph{read}.  A selected point from the jar, quadrature as bracket source, dyadic ownership, process identity from
equal count, laboratory containment or proximity, physical coupling, and every world answer absent apparatus remain
\emph{reserved}.
\section{Nodes, quadrature, and root finding}

Quadrature, root finding, and continued-fraction enclosure answer different questions.  A quadrature rule approximates
an integral.  A root algorithm seeks or encloses a zero of a declared function.  A continued-fraction report brackets
the internal ratio through rational walls.  Shared nodes, counts, or numerals do not merge these reports.

The three-point Gauss--Legendre rule on \([-1,1]\) uses
\[
\xi_1=-\sqrt{3/5},\qquad \xi_2=0,\qquad \xi_3=\sqrt{3/5},
\]
and weights
\[
                         w_1=5/9,\qquad w_2=8/9,\qquad w_3=5/9.
\]
For an interval \([a,b]\), map
\[
x_j=\frac{a+b}{2}+\frac{b-a}{2}\xi_j
\]
and compute the finite weighted sum
\[
Q_3(f)=\frac{b-a}{2}\sum_{j=1}^3w_jf(x_j).
\]
This is a marked diagnostic.  It is exact for polynomials through degree five under exact nodes, weights, and
arithmetic.  For general \(f\), an error claim requires smoothness and a demonstrated derivative or analytic bound.

The irrational nodes must be approximated or enclosed.  One robust method brackets \(z=\sqrt{3/5}\) by positive
rationals \(\ell<u\) satisfying
\[
                         5\ell^2<3<5u^2,
\]
then refines by bisection or a safeguarded rational update.  The node certificate retains the exact inequalities,
iteration trace, stop width, and propagation of node uncertainty through \(f\) and the weighted sum.  A decimal node
without an enclosure supplies only a formatted approximation.

\paragraph{Worked example 1: polynomial quadrature (computation).}
For \(f(x)=x^4\) on \([-1,1]\),
\[
Q_3(f)=\frac{5}{9}\left(\frac35\right)^2+
\frac89(0)+\frac59\left(\frac35\right)^2=\frac25,
\]
which equals the exact integral.  This validates the declared rule on this input.  It does not make every three-point
sum exact.

Root finding begins with a function \(g:D\to\mathbb R\) or an exact rational predicate derived from it.  A sign-change
bracket \([\ell,u]\subset D\) requires
\[
                         g(\ell)g(u)<0.
\]
Existence of an interior root additionally uses continuity on \([\ell,u]\).  Uniqueness requires a separate condition,
such as a derivative of fixed nonzero sign.  Endpoint signs, continuity, and uniqueness are distinct certificate rows.

A safeguarded iteration proposes a Newton, secant, or interpolation step only when it lies in the current bracket and
passes domain checks; otherwise it takes the midpoint.  After evaluating the new point, it retains the subinterval
with opposed endpoint signs.  The trace records proposals, fallbacks, residuals, widths, accepted slices, rejected
slices, and resource use.  A small residual without conditioning information is not a root-error bound.

\paragraph{Worked example 2: exact crossing enclosure (computation).}
For
\[
                         g(d)=5d^2-18,
\]
the rational points \(17/9\) and \(2\) satisfy
\[
                         g(17/9)=-13/81<0,\qquad g(2)=2>0.
\]
Continuity gives a root in the open distance interval \((17/9,2)\); positivity and strict increase for \(d>0\)
identify the positive root as \(\sqrt{18/5}\).  Mapping these walls through the declared internal inverse formula is
the continued-fraction jar construction, not a quadrature result.

\paragraph{Counterexample 3: residual alone is not enclosure (computation).}
A point \(x\) with small \(|g(x)|\) may be far from a root when \(g\) is flat, badly scaled, or contaminated by
rounding.  Without a sign-change bracket or a demonstrated inverse-slope bound, the report owns a residual only.

Refinement studies require one compatible asymptotic family.  For nested quadrature partitions or root discretizations
with
\[
                         A(h)=A^\ast+Ch^p+O(h^{p+1}),
\]
Richardson extrapolation may cancel the leading term when \(p\) is derived or independently estimated from at least
three compatible levels, raw differences decrease, observed order stabilizes, and a demonstrated remainder bound
holds.  Direct and stepped internal routes are not automatically such levels, and choosing \(p\) to hit a target is
rejected.

\paragraph{Counterexample 4: three nodes are not resolution-depth three (computation).}
The quadrature rule counts nodes.  The internal jar's \(K=3\) counts whole continued-fraction partial quotients.
Quadrature nodes and weights neither generate the convergents nor define the jar width.  The jar is not a quadrature
remainder.

The external laboratory inverse alpha may be reported as \(137.036\) only as a separately sourced physical
observation summary.  It belongs to an apparatus, physical model, units and conventions, calibration chain,
uncertainty, data reduction, and provenance.  The covariant acquisition is
\[
                         (N,s)\longmapsto(N',s',y,\Gamma),
\]
and retains sensor loading, actuation, energy and material exchange, latency, hysteresis, drift, saturation, geometry,
environment, prior state, rejected outcomes, back-action, and alternative causal models.  The world answers this
intervention; the number is not produced by the internal model.

Internal \(\alpha_c\), its direct and stepped diagnostics, its count-to-three open jar, and the laboratory report are
different typed objects.  This section makes no nearness, containment, prediction, capture, or real-versus-model
staging.  A shared symbol is forbidden until typed conversion, apparatus calibration, uncertainty, and provenance
justify a particular comparison.

The comparison pipeline reports:
\[
\begin{array}{ll}
\text{exact equality}&\text{equality after typed conversion},\\
\text{resolution equality}&\text{equal labels under one declared quantizer},\\
\text{uncertainty compatibility}&\text{overlap of declared regions}.
\end{array}
\]
Each independently passes, fails, or remains unresolved.  Same-process requires matching inputs, numerical
transitions, traces, costs, models, apparatus action, configuration, and provenance.  A combined agreement bit may
not erase which relation was tested.

Budgets distinguish node-enclosure iterations, function evaluations, weighted arithmetic, root proposals and
fallbacks, bracket refinements, exact and floating growth, derivative or continuity checks, asymptotic levels,
uncertainty propagation, apparatus samples, and replay.  Stops include \textsc{complete},
\textsc{node-domain}, \textsc{node-enclosure-budget}, \textsc{function-domain},
\textsc{nonfinite-evaluation}, \textsc{quadrature-budget}, \textsc{error-bound-unavailable},
\textsc{no-sign-change}, \textsc{continuity-unestablished}, \textsc{root-budget},
\textsc{uniqueness-unresolved}, \textsc{refinement-incompatible}, \textsc{order-unstable},
\textsc{uncertainty-unresolved}, \textsc{calibration-missing}, and \textsc{provenance-incomplete}.

A certificate records diagnostic type, function and domain, exact or enclosed nodes, weights, interval transform,
consumed evaluations, weighted products and sum, quadrature error assumptions, root endpoints and signs, continuity
and uniqueness conditions, proposals and fallbacks, residuals, accepted and rejected slices, refinement family and
observed order, arithmetic and uncertainty bounds, budgets, stops, checked scope, unresolved frontier, releases,
snapshot hash, configuration, protocol, run identity, platform, deterministic inputs, declared nondeterministic
inputs, and---for physical comparison---the complete apparatus interaction record.

Governance fixes diagnostic type, domains, node and weight conventions, approximation and enclosure methods,
quadrature transform, root predicate, continuity and uniqueness evidence, iteration and fallback rules, refinement
and error model, physical comparison source, apparatus requirements, uncertainty, budgets, stop priority, comparison
relations, and certificate grammar.  A changed field begins a new report.

The action ledger is exact.  Diagnostic type, function, domains, nodes, weights, algorithms, smoothness and error
models, physical apparatus model, comparisons, budgets, and certificate grammar are \emph{posited}.  Polynomial
exactness, node-enclosure implications, weighted-sum identities, conditional root existence and uniqueness,
safeguarded-bracket invariants, conditional Richardson cancellation, apparatus calibration implications, and
certificate implications are \emph{demonstrated}.  Evaluations, node bounds, sums, signs, residuals, brackets,
iterations, errors, physical observations, stops, and provenance are \emph{read}.  General exactness from three
nodes, a root from residual alone, quadrature as continued-fraction source, target-fitted extrapolation, laboratory
capture or containment, identity of \(\alpha_c\) with laboratory alpha, and every physical bridge without apparatus
remain \emph{reserved}.
